% Excerpt from arXiv:2404.19737, "Better & Faster Large Language Models via Multi-token Prediction", CC BY 4.0.
% https://arxiv.org/abs/2404.19737
\documentclass{article}

\begin{document}

\subsection{Searching for the optimal $n$}
\label{sect:optimal-future}

To better understand the effect of the number of predicted tokens, we did comprehensive ablations on models of scale 7B trained on 200B tokens of code. We try $n=1,2,4,6$ and $8$ in this setting. Results in table \ref{tab:future_tokens_results} show that training with 4-future tokens outperforms all the other models consistently throughout HumanEval and MBPP for pass at 1, 10 and 100 metrics: +3.8\%, +2.1\% and +3.2\% for MBPP and +1.2\%, +3.7\% and +4.1\% for HumanEval. Interestingly, for APPS/Intro, $n=6$ takes the lead with +0.7\%, +3.0\% and +5.3\%. It is very likely that the optimal window size depends on input data distribution. As for the byte level models the optimal window size is more consistent (8 bytes) across these benchmarks. 


\begin{figure}[t!]
    \centering
    \includegraphics[width=0.8\linewidth]{img/induction.pdf}
    \caption{\textbf{Induction capability of $n$-token prediction models.} 
     Shown is accuracy on the second token of two token names that have already been mentioned previously. Shown are numbers for models trained with a next-token and a 2-token prediction loss, respectively, with two independent runs each. The lines denote per-loss averages. For small model sizes, next-token prediction models learn practically no or significantly worse induction capability than 2-token prediction models, with their disadvantage disappearing at the size of 100M nonembedding parameters.
}

    \label{fig:induction_size}
\end{figure}

\begin{figure}[t!]
    \centering
    \includegraphics[width=\linewidth]{img/backward_order.png}
    \caption{\textbf{Order of the forward/backward in an $n$-token prediction model with $n = 2$ heads.} By performing the forward/backward on the heads in sequential order, we avoid materializing all unembedding layer gradients in memory simultaneously and reduce peak GPU memory usage.}
    \label{fig:backward}
\end{figure}

\end{document}
